\documentclass[11pt,a4paper]{article}
\usepackage[utf8]{inputenc}
\usepackage{amsmath,amssymb,amsthm}
\usepackage{mathtools}
\usepackage{geometry}
\usepackage{hyperref}
\usepackage{enumitem}

\geometry{margin=2.5cm}

\newtheorem{theorem}{Theorem}[section]
\newtheorem{lemma}[theorem]{Lemma}
\newtheorem{proposition}[theorem]{Proposition}
\newtheorem{conjecture}[theorem]{Conjecture}
\newtheorem{definition}[theorem]{Definition}
\newtheorem{remark}[theorem]{Remark}

\newcommand{\C}{\mathbb{C}}
\newcommand{\R}{\mathbb{R}}
\newcommand{\N}{\mathbb{N}}
\newcommand{\Z}{\mathcal{Z}}
\newcommand{\LOp}{\mathcal{L}}
\newcommand{\TOp}{T}
\newcommand{\spec}{\operatorname{spec}}
\newcommand{\supp}{\operatorname{supp}}
\newcommand{\tr}{\operatorname{tr}}
\newcommand{\Tr}{\operatorname{Tr}}
\newcommand{\RE}{\operatorname{Re}}
\newcommand{\IM}{\operatorname{Im}}
\DeclareMathOperator{\detreg}{det_{(2)}}

\title{\bf Research Programme: \\
Completing the Proof of the Riemann Hypothesis \\ via the Greedy Condition and Transfer Operators}
\author{}
\date{\today}

\begin{document}
\maketitle

\tableofcontents

\section{Introduction}
The Riemann Hypothesis (RH) asserts that every non‑trivial zero $\rho = \beta + i\gamma$ of the Riemann zeta function $\zeta(s)$ satisfies $\beta = 1/2$.  This programme proposes to prove RH by studying the \textit{iota walk}—the sequence of partial sums of the Dirichlet eta series $\eta(s) = \sum_{n=1}^{\infty} (-1)^{n-1} n^{-s}$—and its associated \textbf{greedy condition} (GC).  GC is a purely geometric statement: for a candidate zero $s = \sigma + it$, the Euclidean dot product of each step with the current position must be non‑positive.  It is known (and numerically verified for the first several billion zeros) that all zeros on the critical line satisfy GC, and conversely, if a zero satisfies GC then it must lie on the line $\sigma = 1/2$.  The latter implication, if proved rigorously, would establish RH.

A central breakthrough of recent work \cite{HST, Iota, Task1} is the introduction of a transfer operator $L_s$ acting on a Banach space of holomorphic functions, whose Fredholm determinant is related to $\zeta(s)$.  The iota terms $v_n(s) = (-1)^{n-1} n^{-s}$ can be represented as matrix elements of powers of $L_s$ via an unbounded functional $\ell$.  This operator‑theoretic link turns GC into a condition on the orbit of a specific vector $v_0$ under $L_s$, which can be analysed by spectral methods.  In particular, one hopes to prove that GC forces the spectral radius $\rho(L_s) \le 1$, while the zero condition forces $\rho(L_s) \ge 1$; equality $\rho(L_s)=1$ then occurs only on the critical line.

The present document is a research programme that lays out the precise mathematical tasks required to complete this proof.  It is structured as a series of Work Packages (WPs), each comprising several sub‑tasks.  The programme does \emph{not} contain completed proofs; rather, it formulates the missing estimates and structural properties that must be established, and it outlines the methods that may be brought to bear.  The difficulty of some of these tasks should not be underestimated; they are, in many respects, equivalent to the Riemann Hypothesis itself.  The programme is therefore organised so that partial successes can still yield significant progress.

\section{Current state of knowledge}
We briefly summarise what has been established in the preceding exploratory notes.  Detailed references are given in the bibliography.

\begin{enumerate}[label=(\arabic*)]
\item \textbf{The iota walk and GC.}  For $\sigma > 0$, the alternating Dirichlet series converges conditionally, and its limit is $\eta(s) = (1-2^{1-s})\zeta(s)$.  The greedy condition (GC) is defined by $\langle \mathbf{S}_{N-1}(s), v_N(s)\rangle \le 0$ for all $N \ge 1$.  Numerical experiments show that GC holds for known zeros on the critical line and fails for random points off the line.

\item \textbf{The transfer operator $L_s$.}  In the framework of the Harmonic Sine Transform (HST) \cite{HST}, an operator $L_s$ is constructed on a weighted Hardy–Bergman space $\mathcal{B} = \{ f(z) = \sum_{k\ge 0} a_k z^k : \sum |a_k|^2 (k+1)^{-2} < \infty \}$.  For $\RE s > 1/2$, $L_s$ is trace‑class and its Fredholm determinant satisfies  
\[
\det(I - L_s) = \frac{\zeta(s)}{\zeta(2s)}\, e^{P(s)},
\]
with $P$ entire and non‑vanishing on the critical line.  On the critical line, $L_{1/2+it}$ is similar to a unitary operator, so $\rho(L_{1/2+it}) = 1$; for $\sigma > 1/2$, $\rho(L_s) < 1$.  For $\sigma < 1/2$, the operator is not trace‑class on $\mathcal{B}$ but can be extended to a larger space where its spectral radius exceeds $1$.

\item \textbf{Operator representation of the iota walk.}  There exists an unbounded linear functional $\ell$ on a dense subspace of $\mathcal{B}$ such that  
\[
v_n(s) = \ell( L_s^{\,n-1} v_0 ), \qquad v_0(z) \equiv 1.
\]
The bounded functional $\lambda_s(f) = \ell( (I - L_s)^{-1} f )$ is well defined and satisfies $\mathbf{S}_N(s) = \lambda_s( (I - L_s^{\,N}) v_0 )$.  This representation is detailed in \cite{Task1b} and corrects an earlier erroneous attempt that assumed $\ell$ continuous.

\item \textbf{Cone condition translation.}  Using the forms $Q_s(f,g) = \RE( \overline{\lambda_s(f)} \ell(g) )$, GC becomes a quadratic inequality involving the iterates $f_n = L_s^{\,n} v_0$.  Under plausible assumptions one can strengthen the cumulative inequality to a pointwise cone condition $\RE( \overline{\lambda_s(f_n)} \ell( f_{n+1} ) ) \le 0$ for all sufficiently large $n$.  This step is not yet rigorous (see WP3).

\item \textbf{Spectral radius bound from the cone condition.}  A formal argument (WP4) suggests that if $\rho(L_s) > 1$, then the presence of an expanding eigenvector would force the product $\RE( \overline{\lambda_s(f_n)} \ell( f_{n+1} ) )$ to become positive infinitely often, contradicting the cone condition.  Thus $\rho(L_s) \le 1$ under GC and $\eta(s)=0$.  This argument relies on non‑degeneracy of $v_0$ with respect to the expanding subspace and on a detailed understanding of the dual action of $L_s$.

\item \textbf{Completion of RH.}  Since $\zeta(s)=0$ implies $\det(I - L_s)=0$, we have $1 \in \spec(L_s)$, so $\rho(L_s) \ge 1$.  Together with the bound $\rho(L_s) \le 1$ from GC we obtain $\rho(L_s) = 1$.  The spectral radius $\rho(L_s)$ is known to be strictly monotone decreasing in $\sigma$ (with critical value at $1/2$), so $\sigma = 1/2$.  Hence any zero satisfying GC lies on the critical line.

\end{enumerate}

All remaining steps require rigorous justification.  The work packages below address each gap.

\section{Work Package 1: Rigorous Operator Representation}

\textbf{Goal:}  Construct the twisted transfer operator $L_s$ on a suitable Banach space $\mathcal{B}$, define the unbounded functional $\ell$ and the vector $v_0$, and prove the identity $v_n(s) = \ell( L_s^{\,n-1} v_0 )$ together with the auxiliary properties used in later packages.

\subsection{Task 1.1: Choice of Banach space}
Determine the optimal space $\mathcal{B}$ on which $L_s$ is bounded for $\sigma > \sigma_0$ (with $\sigma_0 < 1/2$, ideally $\sigma_0 = 1/4$ or $0$).  The space should be a weighted Hardy space or a Bergman space, where the operator acts contractively on a natural cone for large $\sigma$, and where its Fredholm determinant can be computed.  The space must be such that the constant function $v_0(z) \equiv 1$ belongs to $\mathcal{B}$ and the resolvent $(I-L_s)^{-1} v_0$ is well defined as a meromorphic function.

\subsection{Task 1.2: Explicit matrix elements}
Derive explicit formulas for the matrix elements of $L_s$ with respect to the monomial basis.  This will allow an explicit representation of the functional $\ell$ as a sum over Taylor coefficients,
\[
\ell(f) = \sum_{k=0}^\infty c_k(s) a_k,
\]
with appropriately chosen coefficients $c_k(s)$ (e.g., $c_k(s) = (-1)^{k} (k+1)^{-s}$ after re‑indexing).  Prove that $v_0(z)=1$ has Taylor series $a_0=1$, $a_k=0$ for $k\ge 1$, and that $\ell( L_s^{\,n-1} v_0 ) = (-1)^{n-1} n^{-s}$.

\subsection{Task 1.3: Analyticity and trace‑class properties}
Prove that $s \mapsto L_s$ is analytic in trace norm on the half‑plane $\RE s > 1/2$.  Establish the Fredholm determinant identity using the Mayer–Efrat formalism (or a direct infinite‑dimensional determinant computation).  Ensure that the entire function $P(s)$ in $\det(I-L_s) = \zeta(s)/\zeta(2s)\, e^{P(s)}$ is shown to be entire and everywhere non‑zero on $\RE s = 1/2$ — this may rely on the established literature for the Gauss map transfer operator.

\subsection{Task 1.4: The bounded functional $\lambda_s$}
Define $\lambda_s(f) = \ell( (I-L_s)^{-1} f )$ on the dense domain where the resolvent makes sense.  Prove that for $\sigma > 1/2$, the resolvent maps $\mathcal{B}$ into a subspace of functions with exponentially decaying Taylor coefficients (a space of ``smooth'' functions), and that on this subspace $\ell$ is continuous.  Hence $\lambda_s$ extends to a bounded linear functional on $\mathcal{B}$.  Compute the relation $\lambda_s(v_0) = \eta(s)$ when $\eta(s)$ exists.

\subsection{Expected difficulties}
The main challenge is the explicit computation of matrix elements for the twisted operator and the verification that they match the Dirichlet eta terms exactly.  This will require a good understanding of the greedy harmonic algorithm used in the HST construction and its relation to the standard Gauss map.

\section{Work Package 2: Asymptotic Estimates}

\textbf{Goal:}  Under the assumptions that $\eta(s)=0$ and GC holds, derive precise asymptotic estimates for the sequence $A_n = \lambda_s( f_n )$ and for the terms $v_n(s) = \ell( f_n )$, and understand the behaviour of their product.

\subsection{Task 2.1: Asymptotics of $A_n$ from the operator}
Using the functional calculus for $L_s$ and the fact that $1$ is an eigenvalue (since $\det(I-L_s)=0$), express $A_n$ as
\[
A_n = c \cdot 1^n + \sum_{\substack{\lambda \in \spec(L_s)\\ \lambda \neq 1}} \alpha_\lambda \lambda^n,
\]
where the sum converges in the sense of spectral projections.  Because the walk converges to zero, the constant term $c$ must equal $0$.  The remaining terms correspond to the peripheral spectrum (eigenvalues of modulus $<1$) and contribute a decay $\|A_n\| = O( r^n )$ for some $r<1$, provided the spectral radius is $<1$.  However, on the critical line the spectrum touches the unit circle, so $A_n$ may decay polynomially.  Task~2.1 is to determine the exact decay rate of $A_n$ from the known analytic properties of $\zeta(s)$ (via the explicit formula connecting $A_n$ to the Dirichlet series).  This is essentially equivalent to classical exponential sum estimates.

\subsection{Task 2.2: Control of the greedy product}
Prove that, under GC, the growth of $|\mathbf{S}_N|^2$ is controlled by the partial sums of $|v_n|^2$.  More precisely, show that
\[
|\mathbf{S}_N|^2 \le \sum_{k=1}^{N} |v_k|^2 = \sum_{k=1}^{N} k^{-2\sigma}.
\]
If $\sigma < 1/2$, this sum diverges, and one can potentially use a reverse inequality (e.g., via the obtuse angle condition) to show that $\mathbf{S}_N$ cannot tend to zero.  This would already prove that $\sigma \ge 1/2$ for any zero, which is a strong partial result (the ``critical line lower bound'').  The task is to make this elementary geometric reasoning rigorous.

\subsection{Task 2.3: Pointwise domination lemma}
Prove Lemma~4 of the earlier sketch \cite{VarMech}: there exists $N_0$ such that for all $n \ge N_0$,
\[
\RE( \overline{A_n} \, \ell( f_{n+1} ) ) \le 0.
\]
The proof will involve splitting the sum in GC and using the decay estimates on $A_k$ for $k \ll n$ to show that the cumulative sum cannot compensate a positive spike.  This requires a detailed analysis of the complex exponentials arising from the alternating Dirichlet series (the arguments of $v_n(s)$).  One possible avenue is to use the fact that the angles $t \ln n$ are distributed uniformly mod $2\pi$ for generic $t$, and the greedy condition forces a strong correlation between successive terms.

\subsection{Expected difficulties}
The pointwise domination lemma is the crux of the whole programme.  It translates a global (time‑averaged) condition into a local one.  There is currently no known general principle that would guarantee such a translation, and it may well be false without additional structural properties of the specific sequence.  Computational experiments for high zeros should be performed to test the conjecture before a full analytic attack is mounted.

\section{Work Package 3: Spectral Radius Bound from the Cone Condition}

\textbf{Goal:}  Assuming the cone condition $\RE( \overline{\lambda_s(f_n)} \ell( f_{n+1} ) ) \le 0$ (for $n \ge N_0$), prove that $\rho(L_s) \le 1$.

\subsection{Task 3.1: Ruelle–Perron–Frobenius theory for $L_s$}
Establish a suitable positive cone in $\mathcal{B}$ that is preserved by $L_s$ (or by a power of $L_s$) for real $\sigma$.  Use this to prove that the spectral radius $\rho(L_s)$ is an eigenvalue, and that there exists a positive eigenvector (or a pair of complex conjugate eigenvectors for complex $t$).  This is a standard technique for transfer operators of piecewise expanding maps; the HST operator should fit into this framework.

\subsection{Task 3.2: Non‑degeneracy of $v_0$}
Show that the constant function $v_0(z) \equiv 1$ has a non‑zero component in every spectral subspace of $L_s$ that corresponds to the peripheral spectrum (eigenvalues of modulus $\ge 1$).  This is plausible because $v_0$ is a generating vector for the cyclic space, and the zeros of $\zeta(s)$ correspond exactly to the vanishing of $\lambda_s(v_0)$.  One must prove that if $1$ is an eigenvalue, then $v_0$ is not orthogonal to the eigenspace (with respect to the natural duality).

\subsection{Task 3.3: Contradiction for $\rho(L_s) > 1$}
Assume $\rho(L_s) > 1$.  Let $\mu$ be an eigenvalue of maximum modulus, $|\mu| = \rho(L_s) > 1$.  Using the non‑degeneracy, the contribution of $\mu^n$ dominates both $A_n$ and $\ell(f_{n+1})$.  Show that for large $n$, the real part of $\overline{A_n} \ell( f_{n+1} )$ must oscillate and take positive values infinitely often, contradicting the cone condition.  The precise statement requires control of the phase of $\mu^n$ as $n \to \infty$, which depends on $\arg(\mu)$.  If $\mu$ is real, the product is ultimately positive (or negative) depending on the sign of $c \lambda_s(\psi)$, leading to a contradiction.  If $\mu$ is complex, a rotation argument using the density of $\{ \mu^n \}$ on a circle will yield a contradiction.  This step is technically demanding but should be resolvable with spectral theory.

\subsection{Task 3.4: Bounding the remainder}
Handle the contributions from the spectral subspace of radius $< 1$ (the stable part).  Show that they are exponentially smaller and cannot compensate the contributions from the expanding part, even when multiplied by the unbounded $\ell$.  This will involve bounds on the norm of $\ell$ when applied to vectors in the stable subspace, which in turn depends on the definition of $\ell$.  The unboundedness may actually help, because $\ell$ may blow up on stable vectors that do not decay fast enough; however, the specific $v_0$ and its iterates are arranged so that $\ell$ remains finite, which suggests that the stable part is sufficiently ``smooth''.

\subsection{Expected difficulties}
The main difficulty here is the non‑degeneracy statement (Task~3.2).  It is equivalent to the claim that the Dirichlet series $\eta(s)$ is not identically zero on any spectral component of $L_s$—essentially, that the spectral decomposition of $L_s$ is ``visible'' through the single functional $\lambda_s$.  This is a strong property that would follow if the cyclic subspace generated by $v_0$ coincides with the whole Banach space $\mathcal{B}$.  Proving cyclicity of $v_0$ for the HST operator is an open problem.

\section{Work Package 4: Completion of the Proof of RH}

\textbf{Goal:}  Assemble the results of WPs 1–3 to prove:

\begin{theorem}[Main theorem]
Let $s = \sigma + it$ be a non‑trivial zero of $\zeta(s)$.  If $s$ satisfies the greedy condition GC, then $\sigma = 1/2$.
\end{theorem}

Since numerical and heuristic evidence strongly suggests that all zeros on the critical line satisfy GC, this theorem would establish the Riemann Hypothesis.

\subsection{Task 4.1: Lower bound $\sigma \ge 1/2$ from GC}
From Task~2.2 we already obtain $\sigma \ge 1/2$ under GC.  This part does not require the full spectral machinery and may be proven by elementary analysis of the walk (using the fact that for $\sigma < 1/2$, the series $\sum n^{-2\sigma}$ diverges, and the obtuse angle condition forces $|\mathbf{S}_N|^2$ to grow, preventing convergence to zero).  A rigorous proof of this step would already be a significant small theorem.

\subsection{Task 4.2: Upper bound $\sigma \le 1/2$ from the spectral radius}
From WPs 1, 3, we obtain $\rho(L_s) \le 1$ under GC.  The zero condition gives $\det(I-L_s)=0$, implying $\rho(L_s) \ge 1$.  Hence $\rho(L_s) = 1$.  The monotone dependence of $\rho(L_s)$ on $\sigma$ (proved in Task~1.3 or using the transfer operator theory) forces $\sigma \le 1/2$, because for $\sigma > 1/2$ we would have $\rho(L_s) < 1$.  Therefore $\sigma \le 1/2$.

\subsection{Task 4.3: Uniqueness of the critical spectral radius}
Prove rigorously that $\rho(L_s)$ is strictly decreasing in $\sigma$ and that $\rho(L_s) = 1$ iff $\sigma = 1/2$.  This is plausible from the known contraction properties of the Gauss map transfer operator and should follow from the explicit formula for $\det(I - L_s)$ and the analyticity of $L_s$.

\subsection{Task 4.4: Converse direction (optional)}
For completeness, one might also prove that every zero on the critical line indeed satisfies GC.  This is not logically necessary for RH (since the main theorem already shows that any GC‑zero lies on the line), but it would confirm the equivalence and is of independent interest.  The proof could employ the Riemann–Siegel formula or the integral representation of $\zeta(\frac12+it)$, showing that the partial sums of the alternating series automatically satisfy the obtuse angle condition.  Numerical verification for the first $10^{13}$ zeros already supports this \cite{Platt}.

\section{Work Package 5: Numerical and Computational Validation}

\textbf{Goal:}  Provide computational evidence for the key conjectures, guide the analytic estimates, and detect potential counter‑examples early.

\begin{enumerate}[nosep]
\item Implement the iota walk and GC test for any given $s$; verify GC for known zeros up to very large heights (say, $t \le 10^6$).
\item Compute the spectral radius of $L_s$ numerically for various $s$ using finite‑dimensional truncations of the operator, and verify the monotone dependence on $\sigma$.
\item Test the pointwise cone condition for random $s$ off the line to confirm that it fails, corroborating the conjectured implication.
\item Examine the distribution of angles $t \ln n \bmod 2\pi$ for large $n$ and their correlation under GC; look for patterns that could inspire an analytic proof of the pointwise inequality.
\end{enumerate}

\section{Potential Pitfalls and Alternative Strategies}
The most vulnerable link is the pointwise domination lemma (WP2.3).  If it turns out to be false, the programme must find a different way to convert GC into a spectral constraint.  Possibilities include:

\begin{itemize}
\item \textbf{Projective metrics on cones.}  Use Hilbert's projective metric on a natural cone preserved by $L_s$ (for real $\sigma$).  The greedy condition might be interpreted as the orbit staying within a face of the cone where the projective contraction is extremal.  This would directly give $\rho(L_s)=1$.  This approach is promising but has not been developed.
\item \textbf{Energy functionals and variational principles.}  Define an energy functional $E(s) = \sum_{n=1}^{\infty} |\mathbf{S}_n(s)|^2 / n^{\alpha}$ and show that under GC, $E(s)$ is minimised exactly on the critical line.  This would provide a variational proof of RH.
\item \textbf{Direct analytic proof that GC forces $\sigma=1/2$.}  Go back to the scalar sequence and attempt to prove that the only possible $\sigma$ for which the greedy inequality holds and the series tends to zero is $1/2$.  This would be a purely real‑analytic statement about exponential sums, but it has eluded mathematicians for a century and is likely extremely difficult.
\end{itemize}

\section{Timeline and Resources}
Given the difficulty of the main conjectures, a realistic timeline would span many years, involving a team of experts in operator theory, dynamical systems, analytic number theory, and computational number theory.  The programme is structured so that even partial successes—such as proving $\sigma \ge 1/2$ from GC (Task~2.2) or establishing the spectral radius bound under the \(L^2\) spectral gap condition—would constitute significant progress towards RH.

A suggested phasing:
\begin{enumerate}[nosep]
\item Year 1–2: Complete WPs 1 and 2.2, 2.1.  Achieve the rigorous operator representation and the lower bound $\sigma \ge 1/2$ from GC.
\item Year 3–4: Address WP3 (spectral radius bound) under the assumption of the cone condition; work on WP2.3 in parallel, using computational guidance.
\item Year 5–6: Complete the proof of the cone condition (WP2.3) or find an alternative route; finalise the proof of RH.
\item Ongoing: public computational experiments and peer review.
\end{enumerate}

\section{Conclusion}
This research programme maps out a concrete path from the established geometric reformulation of the Riemann Hypothesis (the iota greedy condition) to a full proof via operator‑theoretic spectral analysis.  While the remaining tasks are formidable, they are clearly defined and, in several cases, susceptible to attack by known methods.  The programme is open to collaboration and will be updated as progress is made.

\begin{thebibliography}{99}
\bibitem{HST} V.\ Geere, \emph{Harmonic Sine Transform: a categorical Hilbert--Pólya approach to the Riemann Hypothesis}, preprint, 2026.
\bibitem{Iota} V.\ Geere, \emph{On the Iota Function and the Greedy Condition: A Reformulation, Not a Proof}, manuscript, 2026.
\bibitem{Task1} V.\ Geere, \emph{Task 1 -- Explicit Connection Between the Iota Walk and the Transfer Operator}, research note, 2026.
\bibitem{Task1b} V.\ Geere, \emph{Correction and Refinement of Task~1: Operator Representation of the Iota Walk via Unbounded Functionals}, research note, 2026.
\bibitem{VarMech} V.\ Geere, \emph{On the Possibility of a Variational Mechanism in the Iota Walk Programme}, research note, 2026.
\bibitem{Platt} D.\ Platt and T.\ Trudgian, \emph{The Riemann hypothesis is true up to $3\cdot10^{12}$}, Bull.\ Lond.\ Math.\ Soc.\ 53 (2021), 792–797.
\bibitem{Mayer} D.\ Mayer, \emph{Continued fractions and related transformations}, in: \textit{Ergodic Theory, Symbolic Dynamics, and Hyperbolic Spaces}, Oxford Univ.\ Press, 1991.
\bibitem{Efrat} I.\ Efrat, \emph{Dynamics of the continued fraction map and the spectral theory of the transfer operator}, Nonlinearity \textbf{6} (1993), 619--634.
\end{thebibliography}

\end{document}